\documentclass[10pt,twocolumn]{article}

\usepackage[margin=0.72in,columnsep=0.25in]{geometry}
\usepackage[utf8]{inputenc}
\usepackage[T1]{fontenc}
\usepackage{lmodern}
\usepackage{microtype}
\usepackage{amsmath,amssymb,amsthm}
\usepackage{booktabs}
\usepackage{array}
\usepackage{enumitem}
\usepackage{xcolor}
\usepackage{listings}
\usepackage{tikz}
\usetikzlibrary{arrows.meta,positioning,fit,shapes.geometric,backgrounds,calc}
\usepackage[numbers,sort&compress]{natbib}
\usepackage[hyphens]{url}
\usepackage[breaklinks=true,colorlinks=true,linkcolor=blue!60!black,citecolor=blue!60!black,urlcolor=blue!60!black]{hyperref}
\usepackage{doi}

% Breakable monospace for identifiers.
\DeclareUrlCommand{\code}{\urlstyle{tt}}
\expandafter\def\expandafter\UrlBreaks\expandafter{\UrlBreaks\do\-\do\_\do\:}

\theoremstyle{plain}
\newtheorem{theorem}{Theorem}
\newtheorem{proposition}[theorem]{Proposition}
\theoremstyle{definition}
\newtheorem{definition}{Definition}

\newcommand{\sem}[1]{[\![#1]\!]}
\newcommand{\covers}{\sqsupseteq}

\lstset{
  basicstyle=\ttfamily\scriptsize,
  breaklines=true,
  breakatwhitespace=false,
  frame=single,
  framerule=0.4pt,
  rulecolor=\color{black!30},
  backgroundcolor=\color{black!3},
  columns=fullflexible,
  keepspaces=true,
  showstringspaces=false,
  aboveskip=4pt,
  belowskip=2pt
}

\newcolumntype{L}[1]{>{\raggedright\arraybackslash}p{#1}}
\emergencystretch=2em
\setlength{\bibsep}{2pt plus 1pt}

\title{\textbf{Executable Ricardian Mandates: Deterministic, Multi-Anchor Authority Vesting for Autonomous AI Agents}}

\author{Valentin Yanakiev}

\date{November 2026}

\begin{document}
\maketitle

\begin{abstract}
Autonomous agents now merge code, move money and provision infrastructure, but the authority they act under is usually an ambient credential, a stream of \texttt{[y/N]} prompts, or a paragraph in a system prompt. None of these is a document a person can read, sign and later be held to, and none can be enforced without trusting the model. We describe \emph{executable Ricardian mandates}: delegations of authority written in a small controlled English whose canonical text is, byte for byte, both what the principal signs and what the gate evaluates. The text is in bijection with a typed syntax tree, so no compiled policy exists that could say something other than what was read. We argue that this property, more than the signature algorithm, is what gives a signed delegation its evidentiary weight. A delegate may only narrow what it received; the narrowing check is sound, needs no SMT solver, and runs in time quadratic in the size of the two mandates. Quotas are charged atomically to every mandate in a chain, so a budget cannot be multiplied by splitting it among children, and every decision is appended to a hash-chained ledger that can be audited offline. Handfest, a Rust implementation, keeps the evaluator pure, puts signature schemes and storage behind two small traits, and enforces mandates in a Model Context Protocol proxy that sits outside the model.
\end{abstract}

% ============================================================================
\section{Introduction}
\label{sec:intro}

An agent that can call tools can do damage in proportion to the credentials it holds. Today those credentials are usually the user's own: an API token in an environment variable, an SSH key on disk, an OAuth grant scoped to a whole account. Least privilege~\cite{saltzer1975protection} is the first casualty. Once a model has been steered by an injected instruction~\cite{greshake2023indirect}, whatever it does is done with the full authority of the person who launched it, and nothing in the system separates what that person meant from what the model did.

The two usual remedies do not survive real workloads. Confirmation prompts work for the first ten tool calls and are rubber-stamped for the next thousand; they do not work at all for sub-agents that run unattended. Describing the permitted behavior in a system prompt, or asking a second model to judge the first, puts a stochastic component inside the security boundary. Such a decision cannot be replayed, cannot be shown to stay within what the principal allowed, and is not evidence of anything.

What is missing is the instrument organizations have long used to hand someone limited authority: a written mandate. A power of attorney is short. It says who may do what, on whose behalf and within which limits; it is signed; and it can be put in front of a third party. We want the same object for agents, with one more property: a machine must be able to enforce it exactly as written.

The pressure for such an object is not only practical. The EU Artificial Intelligence Act requires high-risk systems to be built so that people can oversee and interrupt them (Art.~14) and to record events automatically (Art.~12)~\cite{aiact2024}, and the revised eIDAS regulation obliges member states to offer wallets with which a natural person can produce qualified electronic signatures~\cite{eidas2024}. Both presuppose something concrete to oversee, to log and to sign.

This paper describes such an object and a system that enforces it. An \emph{executable Ricardian mandate} adapts Grigg's Ricardian contract~\cite{grigg2004ricardian}, a single document that is at once legal prose and machine-readable record and is identified by its hash, from a bilateral contract to a unilateral delegation of authority (the civil-law \emph{mandatum}). Handfest, our implementation, is named after the h{\aa}ndf{\ae}stning, the accession charter in which Danish kings accepted written limits on their power.

The paper makes four contributions.
\begin{itemize}[leftmargin=*,itemsep=1pt,topsep=2pt]
  \item A mandate language whose canonical text is in bijection with its typed syntax tree (Section~\ref{sec:language}), and an argument for why this property is what makes a signed delegation usable as evidence (Section~\ref{sec:wysiwys}).
  \item A narrowing relation for sub-delegation that is sound, is decided in quadratic time without a solver, and is exact on the fragments where exactness is cheap (Section~\ref{sec:narrowing}).
  \item Quota accounting across a whole delegation chain that splitting cannot defeat, with a receipt for every permit and every denial in a hash-chained ledger (Sections~\ref{sec:quotas} and~\ref{sec:ledger}).
  \item An implementation with a pure evaluator, together with five families of attacks on it and on its surroundings, each closed and each covered by a regression test (Sections~\ref{sec:arch} and~\ref{sec:attacks}).
\end{itemize}

We should say plainly what exists. One signature scheme is implemented (Ed25519), one storage backend (SQLite), and two gates: a command-line checker and a proxy for the Model Context Protocol. Passkeys, wallet signatures and qualified signatures fit behind the same trait but have not been written. Revocation is missing and so are benchmarks. Section~\ref{sec:limits} collects these gaps.

% ============================================================================
\section{What the Signature Has to Prove}
\label{sec:wysiwys}

A digital signature proves that the holder of a key applied it to a string of bytes. A dispute about a delegation asks something else: did this person intend to grant this authority? Delegations to agents will be contested in the gap between those two questions, so it is worth being precise about the gap before describing any mechanism.

\paragraph{The law asks about intent.}
Signature statutes are written around intent, not mathematics. In the United States, ESIGN and UETA both define an electronic signature as a sound, symbol or process ``executed or adopted by a person with the intent to sign the record''~\cite{esign2000,ueta1999}, and UETA makes the effect of an attributed signature a question of fact, ``determined from the context and surrounding circumstances at the time of its creation, execution, or adoption''~\cite[\S 9(b)]{ueta1999}. The EU's eIDAS regulation grades signatures by assurance instead: an advanced signature must be ``linked to the data signed therewith in such a way that any subsequent change in the data is detectable'', and a qualified one has ``the equivalent legal effect of a handwritten signature''~\cite[Arts.~25--26]{eidas2014}. None of these texts says that the signatory understood, or even saw, what was signed. That is left to evidence.

\paragraph{Signing what one has not seen.}
Now consider how a delegation to an agent is typically signed. The authority lives in a machine artifact: a JSON scope list, a Rego or Datalog policy, a blob of bytecode, some intermediate representation. The principal is shown something else, such as a summary, a consent screen, or a rendering produced by whatever tool is at hand, and then signs the artifact, or more often its hash. When that delegation is disputed, the signature establishes that a key was applied to a digest. To get from there to intent, the party relying on the delegation has to prove facts about software: which renderer produced the screen, that it rendered this artifact and not another, and that what it displayed means what the artifact does. The principal's answer needs no forensic skill: ``I approved what I was shown. I never saw that file and could not have known what it would do.'' Against a compiled policy this is a credible position, and in practice it is the relying party who must rebut it.

The security literature saw the problem early. Landrock and Pedersen gave the requirement its name, \emph{what you see is what you sign} (WYSIWYS)~\cite{landrock1998wysiwys}, and Weber argued that signatures become binding only on devices with secure input and output, where the signatory can see what is being signed~\cite{weber1998see}. McCullagh and Caelli observed that ``non-repudiation'' means one thing to cryptographers and another to courts: at law a purported signer may always dispute a signature, and no key length changes that~\cite{mccullagh2000nonrepudiation}. A hash over an opaque artifact gives non-repudiation in the first sense only.

\paragraph{One document.}
Grigg's Ricardian contract attacks the problem in the document, not in the device. He defines it as ``a single document'' that is ``easily readable by people (like a contract on paper)'' and ``readable by programs (parsable like a database)'', digitally signed and tied to a unique identifier, the hash of the document~\cite{grigg2004ricardian}. Because there is no second representation, there is nothing for the first to disagree with, and Grigg is explicit about the payoff: with a single signed document and an identifier that links it to every record, ``we can more readily ask the judiciary to accept that the single document which is being presented is indeed the valid contract agreed to by the parties''.

Later systems kept the name and loosened the property. Clack et al.\ build on the Ricardian contract by extracting parameters from legal prose and using them to connect the agreement to standardized code~\cite{clack2016templates}. Accord Project templates bind natural-language text to a data model and put the contract's logic in a separate language, Ergo~\cite{roche2021ergo}; EOSIO's Ricardian contracts are prose templates whose variables are substituted from the data of an action that a WebAssembly contract then executes~\cite{eosio2019ricardian}. In each case the text is shown and the code runs, and only the author's care keeps the two in agreement. We call this the \emph{lying template} problem.

\paragraph{What Handfest signs.}
A Handfest mandate restores the single document and adds a proof obligation. The canonical English text is the only part of a signed envelope that carries authority. The typed tree the gate evaluates is obtained from that text by a parser whose inverse is the canonical renderer (Theorem~\ref{thm:bijection}), and the identifier of the mandate is the SHA-256 digest of the text itself. There is no template, no parameter binding and no attached program. The question ``what did the signer authorize?'' therefore has an answer that needs no expert witness: these sentences. Anyone can print them, and anyone can parse them into the tree the gate used.

This changes what can be repudiated. The signer can still say that the key was stolen, that they were coerced, or that the device lied about what was on the screen. Those defenses exist for ink signatures too, and the strength of the trust anchor, a key file at one end and a qualified signature at the other, decides how far they go. What the signer cannot credibly say is that the signed object means something other than what they read, because the signed object \emph{is} what they read. In the terms of e-signature law, the mandate text is both the record and the evidence of what was adopted. In the terms of security engineering, it gives non-repudiation of content and not merely of key use. A dispute about a mandate becomes a dispute about a dozen lines of constrained English, which courts and auditors already know how to resolve.

\paragraph{The price.}
A language that must be read by its signer and inverted by its parser cannot have loops, user-defined functions, or arithmetic across fields. Sections~\ref{sec:language} and~\ref{sec:narrowing} show that what remains is enough to say who may do what, to which resources, under which conditions and how often, and that the same restriction is what makes narrowing decidable without a solver. WYSIWYS in its original sense also needs an honest display. Plain text keeps the trusted viewer as simple as a viewer can be, but we do not otherwise address the signing device (Section~\ref{sec:limits}).

% ============================================================================
\section{Related Work}
\label{sec:related}

\paragraph{Delegation tokens.}
OAuth~2.0 Token Exchange lets a client trade one token for another that names both a subject and an acting party, through the \texttt{act} claim~\cite{rfc8693}. Every exchange is a round trip to an authorization server, and what comes back is a token whose meaning rests on scope strings defined elsewhere. Macaroons~\cite{birgisson2014macaroons} removed the round trip: the holder of a token can add caveats that restrict it, chained by HMAC, without contacting the issuer, although verification needs the root secret. Biscuit~\cite{biscuit} does the same with public-key signatures and Datalog checks, and UCAN~\cite{ucan} builds public-key delegation chains in which each link may narrow the one before. Handfest takes offline narrowing from this line of work. It differs in what is narrowed, signed sentences in place of a serialized token, and in keeping usage counters for every mandate in a chain, which needs state that a bearer token does not carry.

ERC-5573 (ReCaps)~\cite{erc5573} is the closest relative on the first point. It translates a capability object deterministically into an English statement inside the Sign-In with Ethereum message that the wallet displays, so that the user signs a description of what is delegated. The statement lists the granted abilities for each resource but not the caveats attached to them. The scheme is tied to Ethereum wallet signatures, and it leaves the evaluation of delegations, usage limits and logging to the relying party.

\paragraph{Policy languages.}
Cedar~\cite{cutler2024cedar} is the best current example of an authorization language designed for analysis. Its semantics are modeled in Lean, and its design admits a sound and complete logical encoding, so that an SMT solver such as cvc5~\cite{barbosa2022cvc5} can decide, for example, whether a refactored policy set authorizes the same requests as the original. Cedar policies are written and stored by the party that operates the resource. They are not documents that a delegator signs and hands on. Handfest sits at the other end: each mandate is a signed delegation, and its language is small enough that containment needs no solver. Nothing prevents an organization from running both.

\paragraph{Controlled languages.}
Controlled natural languages show that readable need not mean ambiguous. Attempto Controlled English is a subset of English that translates unambiguously into first-order logic~\cite{fuchs2008ace}, and Catala applies literate programming to statutes, setting each article of law beside the code that formalizes it~\cite{merigoux2021catala}. Both aim at expressive power. We go the other way and keep only what a delegation needs, because the language has to stay invertible and containment has to stay cheap.

\paragraph{Identity.}
Who controls a key is a separate question from what the key signed. Self-certifying identifiers with key event logs~\cite{smith2019keri} and the credentials chained from them~\cite{smith2026acdc}, W3C verifiable credentials~\cite{w3c2025vcdm}, selective-disclosure tokens~\cite{rfc9901,terbu2026sdjwtvc}, the EU wallet architecture~\cite{eudi2026arf} with its trusted lists~\cite{etsi2025tl}, passkeys~\cite{w3c2021webauthn} and typed wallet signatures~\cite{eip712} all answer the first question. Handfest treats them as interchangeable trust anchors. Each would implement one trait (Section~\ref{sec:anchors}) that turns scheme-specific evidence into the same small proof record. Only Ed25519 keys~\cite{rfc8032} are implemented.

\paragraph{Agents, tools and logs.}
The Model Context Protocol (MCP)~\cite{mcp2025spec} has become a common way for agent harnesses to reach tools: JSON-RPC~2.0 messages~\cite{jsonrpc2} over a pipe or over HTTP. Because the harness speaks the protocol and the model does not, a proxy on that channel is a natural place to enforce a mandate (Section~\ref{sec:gates}). Chaining log entries by hash so that later tampering shows goes back at least to Schneier and Kelsey~\cite{schneier1999audit}, and transparency logs such as Certificate Transparency~\cite{rfc6962} add public verifiability. Our ledger uses the first idea and, as Section~\ref{sec:attack-ledger} explains, needs something like the second to resist a writer who can replace the whole file.

% ============================================================================
\section{The Mandate Language}
\label{sec:language}

A mandate is a short text with a fixed order of clauses. Figure~\ref{fig:grammar} gives the grammar and Figure~\ref{fig:example} a delegation in two hops.

\begin{figure}[t]
\begin{lstlisting}
mandate       = "HANDFEST v1" NL
                [ "PARENT " digest NL ]
                "FROM " actor_decl NL
                "TO " actor_decl NL
                "DELEGATION max " uint (" hop"|" hops") NL
                "ALLOW " action_list NL
                "ON " resource_list NL
                [ "WHERE " predicate
                  ( NL "  AND " predicate )* NL ]
                [ "LIMIT " quota
                  ( NL "  AND " quota )* NL ] ;

actor_decl    = actor_kind " " quoted_str
                " (key: " key_id ( ", " attr_kv )* ")" ;
actor_kind    = "human" | "org" | "agent" | "service" ;
attr_kv       = ident ": " scalar ;
action_list   = action_pat ( ", " action_pat )* ;
action_pat    = ident ( "." ident )* [ ".*" ] ;
resource_list = resource_sel ( ", " resource_sel )* ;
resource_sel  = ident " " quoted_glob ;
predicate     = field_path " " (
                  ("=="|">="|">"|"<="|"<") " " scalar
                | "IN {" scalar_list "}"
                | "PREFIX " quoted_str ) ;
quota         = uint ( " call" | " calls" )
              | decimal " " unit_ident
                " sum(" field_path ")" ;
field_path    = ident ( "." ident )* ;
scalar_list   = scalar ( ", " scalar )* ;
scalar        = quoted_str | rfc3339_ts
              | decimal | int | bool ;
bool          = "true" | "false" ;
digest        = "sha256:" hex64 ;
key_id        = ident ":" non_ws ;
NL            = "\n" ;
\end{lstlisting}
\caption{Grammar of the mandate language, from \texttt{specs/grammar.ebnf} (re-wrapped). Lexical rules for identifiers, strings, decimals and timestamps are enforced by the parser and summarized in Definition~\ref{def:mandate}.}
\label{fig:grammar}
\end{figure}

\begin{figure}[t]
\begin{lstlisting}
HANDFEST v1
FROM human "alice" (key: ed25519:9b1e...44a0)
TO agent "dispatcher" (key: ed25519:3c8d...19f2)
DELEGATION max 1 hop
ALLOW git.pr.merge
ON repo "acme/*"
WHERE context.time >= 2026-11-24T09:00:00Z
  AND context.time < 2026-11-24T17:00:00Z
  AND resource.ci_status == "passing"
LIMIT 5 calls
\end{lstlisting}
\begin{lstlisting}
HANDFEST v1
PARENT sha256:8f4e...c901
FROM agent "dispatcher" (key: ed25519:3c8d...19f2)
TO agent "seat-7" (key: ed25519:7a20...e114)
DELEGATION max 0 hops
ALLOW git.pr.merge
ON repo "acme/webshop"
WHERE context.time >= 2026-11-24T10:15:00Z
  AND context.time < 2026-11-24T10:18:00Z
  AND resource.ci_status == "passing"
  AND resource.pr == 492
LIMIT 1 call
\end{lstlisting}
\caption{A root mandate signed by a person (top) and a child that the dispatcher signs for one of its workers (bottom). The child narrows the repository, the time window and the budget, adds a condition, and forbids further delegation. Keys and the parent digest are abbreviated here; real ones are 64 hexadecimal digits.}
\label{fig:example}
\end{figure}

\begin{definition}[Mandate]
\label{def:mandate}
A mandate is a tuple $m = (\mathit{parent}, \mathit{from}, \mathit{to}, h, A, R, W, Q)$ in which
\begin{itemize}[leftmargin=*,itemsep=0pt,topsep=2pt]
  \item $\mathit{parent}$ is absent or a 256-bit digest;
  \item $\mathit{from}$ and $\mathit{to}$ are actors, each with a kind (human, org, agent or service), a name, a key identifier of the form \emph{scheme}\texttt{:}\emph{value}, and a map of typed attributes;
  \item $h < 2^{32}$ is the number of further delegation hops allowed;
  \item $A$ is a non-empty set of action patterns, each a dotted sequence of identifiers that may end in \texttt{.*};
  \item $R$ is a non-empty set of resource selectors, each a kind and a glob in which \texttt{*} is the only wildcard;
  \item $W$ is a set of predicates, each a field path with a comparison (\texttt{==}, \texttt{>=}, \texttt{>}, \texttt{<=}, \texttt{<}) against a scalar, membership (\texttt{IN}) in a non-empty set of scalars, or a string \texttt{PREFIX};
  \item $Q$ is a set of quotas, each either $n$ \texttt{calls} with $n < 2^{64}$ or a non-negative decimal bound on the sum of a named field in a named unit.
\end{itemize}
Scalars are strings, 64-bit signed integers, decimals $c \cdot 10^{-s}$ with a 128-bit signed coefficient $c$ and $0 \le s \le 28$, UTC timestamps in whole seconds from year 0000 to year 9999, and booleans. Identifiers match \texttt{[A-Za-z\_][A-Za-z0-9\_-]*}. Strings contain no control characters. Key values consist of visible ASCII characters other than the comma, the parentheses and the double quote.
\end{definition}

The Rust type for mandates has private fields and two constructors, the parser and a builder, both of which validate these rules and then \emph{canonicalize}: clause lists and \texttt{IN} sets are sorted and deduplicated, and decimals are reduced to the shortest form that keeps one fractional digit. Every value of the type is therefore valid and canonical. Write $\mathcal{M}$ for the set of such values.

\subsection{Canonical Text}
\label{sec:bijection}

Two functions connect text and tree. The renderer $\mathsf{render} : \mathcal{M} \to \Sigma^{*}$ prints a mandate in one fixed layout: a line for each clause, predicate and quota, a line feed after every line, the quotation mark and the backslash escaped inside strings, and \texttt{hop} or \texttt{call} in the singular exactly when the count is one. The strict parser $\mathsf{parse}$ is partial. It first reads its input leniently, tolerating extra spaces, CRLF line ends, unsorted or repeated list entries and either grammatical number, and builds a mandate $m$. It then accepts only if $\mathsf{render}(m)$ equals the input byte for byte. Let $\mathcal{C} = \mathsf{render}(\mathcal{M})$ be the set of canonical texts.

\begin{theorem}[Text and tree are in bijection]
\label{thm:bijection}
$\mathsf{parse}$ is defined exactly on $\mathcal{C}$, and $\mathsf{render} : \mathcal{M} \to \mathcal{C}$ and $\mathsf{parse} : \mathcal{C} \to \mathcal{M}$ are inverse to each other.
\end{theorem}

\begin{proof}[Proof sketch]
If $\mathsf{parse}(t) = m$ then $\mathsf{render}(m) = t$ by the final comparison, so the domain of $\mathsf{parse}$ lies inside $\mathcal{C}$ and $\mathsf{render} \circ \mathsf{parse}$ is the identity on it. It remains to show that the lenient reader recovers every $m \in \mathcal{M}$ from $\mathsf{render}(m)$; then $\mathsf{parse}$ is defined on all of $\mathcal{C}$ and $\mathsf{parse} \circ \mathsf{render}$ is the identity. This is a property of the grammar. Strings cannot contain a line feed, so the text splits into lines unambiguously, and each line begins with a fixed keyword in a fixed order. Within a line every token delimits itself: identifiers and field paths are ASCII words; strings are quoted, with the two escapes above; key values exclude the delimiters that can follow them; and the unquoted literals cannot be confused: the booleans are two keywords, a timestamp begins with four digits and a hyphen, a decimal contains a point, and an integer has neither. The reader therefore reconstructs each clause of $m$ exactly, and canonicalizing a tree that is already canonical changes nothing.
\end{proof}

The second half of the argument is by inspection of the grammar and the code. It is exercised by round-trip tests, including the extreme decimal and timestamp values, but it has not been mechanized.

The identifier of a mandate is the digest of its canonical text,
\[
  \mathsf{id}(m) = \mathrm{SHA\text{-}256}\bigl(\mathsf{render}(m)\bigr),
\]
and it is this value that appears in a child's \texttt{PARENT} line, in quota counters and in receipts. By Theorem~\ref{thm:bijection}, an identifier commits to one text and one tree, up to collisions in SHA-256.

% ============================================================================
\section{Narrowing and Quotas}
\label{sec:narrowing}

\subsection{What a Mandate Permits}
\label{sec:semantics}

The gate evaluates a mandate against a \emph{context} $e = (a, k, r, F, \mathit{now}, u)$. Here $a$ is a concrete action, a dotted sequence of identifiers; $k$ and $r$ are the kind and the identifier of a resource, the latter a string without control characters; $F$ is a finite map from field paths to scalars, the \emph{facts}; $\mathit{now}$ is a timestamp; and $u$ is the usage recorded so far for the mandate being evaluated, namely a number of permitted calls and, for each pair of unit and field, a non-negative decimal total. The value of a field in $e$ is the timestamp $\mathit{now}$ for the reserved path \texttt{context.time} and the fact $F(f)$ for every other path $f$.

The policy clauses $(A, R, W, Q)$ of a mandate \emph{permit} $e$ when four conditions hold.
\begin{enumerate}[leftmargin=*,itemsep=0pt,topsep=2pt,label=(\roman*)]
  \item Some pattern in $A$ matches $a$. A pattern without a wildcard matches only itself, and $s\texttt{.*}$ matches every action that extends $s$ by at least one identifier.
  \item Some selector in $R$ has kind $k$ and a glob that matches $r$.
  \item Every predicate in $W$ is satisfied by the value of its field, which must be defined. A comparison holds only between scalars of the same type and never orders booleans; strings are ordered by code point; \texttt{PREFIX} holds only of strings.
  \item Every quota in $Q$ has room. For $n$~\texttt{calls}, the recorded count is below $n$. For a sum quota, the value of its field is a decimal $\delta \ge 0$, and the recorded total plus $\delta$ is a representable decimal no greater than the bound.
\end{enumerate}
We write $\sem{m}$ for the set of contexts that the clauses of $m$ permit. The header of a mandate (parent, actors, hop count) plays no part in $\sem{m}$ and is handled by the chain rules below. A mandate with \texttt{LIMIT 0 calls} is valid and permits nothing.

Ordered by inclusion, the finite intersections of such sets form a meet-semilattice with the empty set, \emph{deny everything}, at the bottom. It has no top, because an \texttt{ALLOW} list is finite and every pattern in it begins with a concrete identifier, so no mandate permits every action. Evaluating a chain of mandates is taking a meet (Section~\ref{sec:quotas}).

\subsection{The Narrowing Relation}
\label{sec:laws}

Say that a pattern $p$ \emph{covers} a pattern $c$ if they are equal, or if $p = s\texttt{.*}$ and $c$ extends $s$, by at least one identifier when $c$ itself has no wildcard. A selector covers another if their kinds are equal and every string matched by the second glob is matched by the first. A set of predicates on one field \emph{implies} a predicate if every scalar that satisfies all of the set satisfies the predicate.

\begin{definition}[Narrowing]
\label{def:narrowing}
A mandate $C$ \emph{narrows} a mandate $P$, written $P \covers C$, when all of the following hold.
\begin{enumerate}[leftmargin=*,itemsep=0pt,topsep=2pt,label=(N\arabic*)]
  \item \emph{Hand-off.} $\mathit{parent}_C = \mathsf{id}(P)$, and the key and the kind of $\mathit{from}_C$ are those of $\mathit{to}_P$.
  \item \emph{Depth.} $h_P \ge 1$ and $h_C \le h_P - 1$.
  \item \emph{Scope.} Every pattern in $A_C$ is covered by a pattern in $A_P$, and every selector in $R_C$ by a selector in $R_P$.
  \item \emph{Conditions.} For every field that $C$ constrains, the predicates of $C$ on that field are satisfiable together. Every predicate of $P$ is implied by the predicates of $C$ on the same field; in particular $C$ constrains that field.
  \item \emph{Quotas.} For every $n$~\texttt{calls} in $Q_P$ there is $n'$~\texttt{calls} in $Q_C$ with $n' \le n$, and for every sum quota in $Q_P$ there is one in $Q_C$ on the same unit and field whose bound is no greater.
\end{enumerate}
\end{definition}

N1 ties a child to one parent text and to the actor that parent addressed. Its kind check stops an agent from describing itself as a human when it delegates onward. N2 bounds the length of a chain by $h_0 + 1$, where $h_0$ is the hop count of the root, and rules out cycles. In Figure~\ref{fig:example} the child satisfies all five conditions.

\begin{theorem}[Narrowing is sound]
\label{thm:sound}
If $P \covers C$ then $\sem{C} \subseteq \sem{P}$.
\end{theorem}

\begin{proof}
Let $e \in \sem{C}$; we check the four conditions for $P$. (i)~Some $c \in A_C$ matches $a$, and some $p \in A_P$ covers $c$. If $p = c$ we are done. Otherwise $p = s\texttt{.*}$ and $c$ extends $s$. If $c$ has no wildcard then $a = c$, which extends $s$ by at least one identifier; if $c$ has one, $a$ extends $c$ by at least one identifier and so extends $s$ by at least one. Either way $p$ matches $a$. (ii)~Some selector in $R_C$ of kind $k$ matches $r$, and a selector in $R_P$ of the same kind matches every string the first one does. (iii)~Let $\varphi$ be a predicate of $P$ on a field $f$. By N4, $C$ has predicates on $f$; since $e \in \sem{C}$ the value $v$ of $f$ is defined and satisfies all of them, so $v$ satisfies $\varphi$. (iv)~For $n$~\texttt{calls} in $Q_P$, take $n' \le n$ from N5: the count in $u$ is below $n'$, hence below $n$. For a sum quota of $P$, take the matching quota of $C$: both read the same fact $\delta$ and the same total in $u$, so the sum is representable and at most the child's bound, which is at most the parent's.
\end{proof}

Theorem~\ref{thm:sound} compares the two mandates under the \emph{same} usage $u$. At run time each mandate has its own counters, and a parent may be nearer exhaustion than its child. The gate therefore never relies on the leaf alone; Section~\ref{sec:quotas} returns to this.

\subsection{Deciding It Without a Solver}
\label{sec:decide}

Containment of policies is hard in general, which is why analyzable languages such as Cedar reduce it to SMT~\cite{cutler2024cedar}. Here the language was cut to fit the check: predicates never relate two fields, and the only disjunction is the finite set of an \texttt{IN}.

\begin{proposition}[Decision procedure]
\label{prop:decide}
The function \texttt{subsumes(P, C)} of \code{handfest-core} succeeds exactly when $P \covers C$. It calls no solver and runs in time $O(n^2)$, where $n$ is the combined length in bytes of the two canonical texts.
\end{proposition}

\begin{proof}[Proof sketch]
N1, N2 and N5 are comparisons of header fields and of quota entries, and pattern covering is a comparison of identifier lists, so only the two semantic notions need an argument.

\emph{Globs.} Collapse runs of \texttt{*} and write the parent glob as literals $p_0, \dots, p_m$ separated by wildcards and the child as $c_0, \dots, c_j$. If $m = 0$ the parent matches one string, and the child is covered iff it is the same literal. If $j = 0$ the child is one string and the parent's matcher is run on it. Otherwise the procedure requires $p_0$ to be a prefix of $c_0$ and $p_m$ a suffix of $c_j$, and then looks for $p_1, \dots, p_{m-1}$ in order, leftmost first and without overlap, inside what remains of $c_0, \dots, c_j$, each $p_i$ within a single $c$. These conditions are sufficient because every string matched by the child contains the child's literals in order. They are necessary because the child matches the string obtained by putting one character that occurs nowhere in the parent glob in place of each wildcard; a match of the parent on that string must place each parent literal inside one child literal, and a leftmost-first search finds such a placement whenever one exists.

\emph{Predicates.} For each field the procedure normalizes the child's predicates into a type, a lower and an upper bound, an optional finite set and an optional prefix. Integers, timestamps and decimals are discrete, so a strict bound becomes a closed one at the neighboring value; decimals are finite in number, and the successor and predecessor of a decimal are found by inspecting its 29 possible scales. For strings, $x > s$ becomes $x \ge s\,\texttt{\textbackslash 0}$, the least string above $s$, and a prefix $q$ becomes the half-open interval from $q$ to the least string above every extension of $q$, or the interval from $q$ upwards when no such string exists. An empty interval, clashing types, an ordering on a boolean, an equality outside the bounds or an \texttt{IN} set with no surviving member is a contradiction, and the child is refused. Otherwise the set of scalars satisfying the child's predicates on the field is exactly the surviving finite set if there is one, and a single value or an interval if not. Each parent predicate is tested against it: member by member for a finite set, and by comparing bounds for an interval. An interval of a discrete type implies a parent \texttt{IN} set iff it has no more values than the set and each is a member. The procedure counts first and enumerates only when the count allows. Decimals fall into 57 contiguous bands according to their sign and the largest scale at which they can be written, and within a band they are equally spaced, so the count is exact when the two ends lie in the same band or in neighboring ones. When a whole band lies between them the interval holds more than $10^{38}$ values, more than any set a mandate can contain, and the answer is no without counting. String intervals written in the language are infinite unless their two ends coincide, so they imply no \texttt{IN} set and no equality.

\emph{Cost.} Pattern and selector lists are compared pairwise, and covering one glob by another takes time at most quadratic in their combined length. Grouping and normalizing the child's predicates takes $O(n \log n)$. A parent predicate is tested against at most the members of one child set, each test being a comparison or a logarithmic set lookup, and an interval is enumerated only up to the size of the parent set it is tested against. Arithmetic on decimals takes a bounded number of steps. Each of these is within $O(n^2)$, and rendering and hashing $P$ are linear.
\end{proof}

As with Theorem~\ref{thm:bijection}, the argument is by case analysis on the code and is backed by tests, not by a mechanized proof.

\paragraph{Where narrowing is stricter than containment.}
The converse of Theorem~\ref{thm:sound} is false on purpose. The clearest case is time. At run time \texttt{context.time} always holds a timestamp, so a parent predicate such as \texttt{context.time >= 0000-01-01T00:00:00Z} excludes nothing. A child that does not mention \texttt{context.time} is refused all the same, because N4 asks the child to constrain every field its parent constrains, and because implication in N4 ranges over all scalars and not only over timestamps. Likewise, a selector must be covered by one parent selector and not by several jointly, and a child with contradictory predicates is refused outright, although a mandate that permits nothing is contained in every other. Each of these refusals costs the delegate a line of text; none can widen authority.

\subsection{Quotas Across a Chain}
\label{sec:quotas}

Narrowing does not bound usage. Take a root mandate with \texttt{LIMIT 5 calls} and one hop of delegation. Its delegate can sign ten children, each a valid narrowing with \texttt{LIMIT 1 call}. If each child were judged alone, ten calls would pass; the same trick multiplies a spending limit. We call this \emph{splitting}.

The gate therefore receives a request together with its whole chain $M_0, \dots, M_k$ and permits it only if two things hold. First, the chain is well formed: $M_0$ has no parent; $M_i \covers M_{i+1}$ for every $i$; each mandate comes with a proof, from the key and the actor kind of its \texttt{FROM} line, whose validity window contains $\mathit{now}$; and the requesting actor's key is the \texttt{TO} key of $M_k$. Second, \emph{every} $M_i$ permits the context built from the request, the facts, $\mathit{now}$ and $M_i$'s own recorded usage. An empty chain is refused. On a permit, and only then, the ledger adds one call to the counter of every mandate in the chain and adds the charged amount to every sum counter that a mandate in the chain declares, in the transaction that also appends the receipt (Section~\ref{sec:ledger}).

\begin{proposition}[Budgets cannot be split]
\label{prop:budget}
Consider a single ledger that only the gate modifies, and let $M$ carry the quota $n$~\texttt{calls}. Then at most $n$ requests are ever permitted on chains that contain $M$. If $M$ carries a sum quota with bound $d$ for a unit and a field, the amounts charged for that unit and field by permitted requests on chains containing $M$ total at most $d$. Every chain that ends in a descendant of $M$ contains $M$, up to collisions in SHA-256.
\end{proposition}

\begin{proof}
The ledger serializes decisions, each in one transaction that reads the counters, evaluates and writes. By induction over the sequence of decisions, the call counter of $M$ equals the number of permits so far on chains containing $M$, since a permit adds one for each mandate of its chain and a denial adds nothing. A permit requires $M$ to permit the context with that counter as usage, that is, the counter is below $n$ beforehand. Hence it never exceeds $n$. The argument for sums is the same, with the exact decimal total in place of the count. For the last claim, a well-formed chain has $\mathit{parent}_{i+1} = \mathsf{id}(M_i)$ and a root without a parent, so a child of $M$ can only occur directly after a mandate whose identifier is $\mathsf{id}(M)$, which is $M$ itself by Theorem~\ref{thm:bijection} unless two texts collide; deeper descendants follow by induction.
\end{proof}

Two remarks keep the proposition honest. It is a statement about one ledger: two ledgers that share no state each start from zero and each will grant the full budget, so a deployment on several hosts needs one shared ledger or a quota service in front of them. And counters are keyed by identifier, that is, by the words. Signing the same text again does not refill a budget, while changing a character makes a new mandate with counters of its own, which for a delegate still charges the unchanged parent.

% ============================================================================
\section{Architecture}
\label{sec:arch}

Handfest is a Rust workspace of five crates with a one-way dependency graph (Figure~\ref{fig:arch}). The evaluator at the bottom is a pure function. Signature schemes and storage sit behind two traits above it, and two front ends, a command-line tool and an MCP proxy, turn the outside world into calls on those three layers. Figure~\ref{fig:seq} follows one delegation from signature to audit.

\subsection{The Pure Core}
\label{sec:core}

The crate \code{handfest-core} contains the mandate type, the parser and renderer, the narrowing check and the evaluator. It is built with \texttt{\#![no\_std]}, uses only the allocator, forbids \texttt{unsafe}, and has one dependency, the \code{sha2} crate. It opens no file, reads no clock and draws no randomness. Everything the evaluator needs arrives as an argument: the chain with its proofs, the request, the facts, the usage snapshot, the timestamp and the hash of the previous receipt.

Its entry point, \texttt{authorize}, returns a verdict and a \emph{decision receipt}. The receipt records the identifiers of the chain, the actor key, the action, the resource, the timestamp, the verdict with its reason, and the hash of the previous receipt; its own hash is SHA-256 over a domain-separated, length-prefixed encoding of those fields. Because the function is pure, running it again on the same inputs reproduces the verdict and the receipt byte for byte. A receipt stores commitments, not inputs: it contains neither the mandate texts nor the facts, so replaying a decision needs those from elsewhere.

Amounts are decimals with a 128-bit coefficient. Addition is carried out on a 256-bit intermediate and fails, without rounding, when the exact result cannot be represented. Floating point appears nowhere in the workspace's handling of amounts, including the JSON reader shared by the two front ends, which keeps number tokens as text (Section~\ref{sec:attack-numeric}).

\subsection{Trust Anchors}
\label{sec:anchors}

A signed mandate travels as a small JSON envelope: the canonical text, an anchor identifier, an assurance level, the signer's actor kind and key, a validity window and a signature. The crate \code{handfest-anchors} defines one trait, \code{AnchorVerifier}, with one method that takes such an envelope and returns either an error or a \emph{verified mandate}: the parsed mandate paired with a proof record of anchor, assurance level, signer key, actor kind and window. The evaluator sees only that record. An \code{AnchorRegistry} picks the verifier by the scheme of the signer's key and then re-checks, for any verifier, that the proof it returned agrees with the envelope and that the mandate it returned renders to exactly the text that was signed. A buggy or hostile verifier cannot substitute another mandate. The proxy verifies through the registry; the command-line tool still calls the Ed25519 verifier directly.

The one verifier that exists is for Ed25519~\cite{rfc8032}. It refuses text that is not canonical, requires the signer to be the mandate's \texttt{FROM} key, and verifies with \code{ed25519-dalek}'s \code{verify_strict}, which is stricter than the RFC in that it also rejects small-order public keys and small-order $R$ components. The signature covers a domain-separated, length-prefixed encoding of every envelope field and not just the text (Section~\ref{sec:attack-envelope}). Assurance levels are an ordered enumeration from local key to qualified signature; the Ed25519 verifier accepts only the lowest, and no policy consumes the level yet.

\subsection{The Receipt Ledger}
\label{sec:ledger}

The crate \code{handfest-ledger} defines the \code{ReceiptLedger} trait: evaluate a request and record the outcome, record a denial that never reached the evaluator, and audit. The implementation is over SQLite, with a table of receipts and two tables of counters.

Evaluation and recording are one \texttt{BEGIN IMMEDIATE} transaction. It reads the hash of the last receipt and the counters of every mandate in the chain, calls \texttt{authorize}, appends the receipt whatever the verdict, updates the counters if the verdict is a permit, and commits. The caller learns the verdict only after the commit. A permitted action that is never carried out therefore still consumes quota, which is the safe direction for the error to go.

Auditing needs no write access. The ledger runs in SQLite's rollback-journal mode, not in write-ahead mode: a write-ahead database can be opened read-only only if its \texttt{-shm} and \texttt{-wal} files already exist, or can be created, or the database is declared immutable~\cite{sqlite_wal}. An auditor who is handed a read-only file in a read-only directory, possibly while a writer is still active, has none of these. An audit opens the file read-only and, inside one read transaction, checks the schema against the expected DDL, runs SQLite's full integrity check, replays the hash chain from the first receipt, compares the number of receipts with the high-water mark SQLite keeps for the receipt sequence, and reconciles every call counter with the permits that name its mandate. Section~\ref{sec:attack-ledger} says what this does and does not detect.

\subsection{Gates}
\label{sec:gates}

\paragraph{Command line.}
The \texttt{handfest} binary has six subcommands: \texttt{fmt} prints the canonical form of a mandate, \texttt{keygen} and \texttt{sign} produce keys and envelopes, \texttt{verify} checks the signatures and the narrowing of a chain, \texttt{check} evaluates one request against a chain and a ledger and exits with status zero only on a permit, and \texttt{audit} verifies a ledger read-only. With \texttt{check}, the chain, the actor key, the facts and the timestamp are all arguments, so whoever invokes it is trusted for them. It is a building block for scripts and hooks, not a boundary against its caller.

\paragraph{MCP proxy.}
The \code{handfest-mcp} binary is the boundary. A harness starts it in place of an MCP server, and it starts the real server as a child process, so that every JSON-RPC message between the two passes through it. The operator who launches it supplies four things the model cannot influence: the signed chain, the key of the actor the session runs as, the ledger, and a \emph{tool policy map}. At start-up the proxy verifies every signature and the narrowing of the chain, and refuses to start the server if either fails.

The policy map is a JSON file that names, for each tool the operator is prepared to expose, the action and resource kind it stands for, where the resource identifier comes from, and which facts accompany a call. A fact is either fixed in the map or lifted, with a declared type, from a named argument of the call itself. A fixed fact always overrides an argument, and \texttt{context.time} comes only from the proxy's clock. Tool names live in this file and not in the code.

On \texttt{tools/call} the proxy maps the call to a request and facts and asks the ledger to evaluate and record it. On a permit the original frame is forwarded unchanged. On a denial the server is never contacted and the client receives a JSON-RPC error with code $-32001$ that carries the reason and the hash of the receipt. A call that cannot be mapped at all, because the tool is unknown or an argument is missing or of the wrong type, is recorded as a denial in the same ledger. On \texttt{tools/list} the proxy removes every tool that is missing from the map or that the chain could not permit, so that the model is not offered what it may not use. Requests other than \texttt{initialize}, \texttt{ping}, \texttt{tools/list} and \texttt{tools/call} are refused, and notifications outside a short allow-list are dropped.

Neither gate pins a root key. Both authenticate each mandate against its own \texttt{FROM} key and trust the launcher to have chosen the right chain. That matches a harness wrapping its own tool servers; a broker serving several principals would also need a list of accepted roots (Section~\ref{sec:limits}).

% ---------------------------------------------------------------------------
% Figure: architecture (full page)
% ---------------------------------------------------------------------------
\begin{figure*}[p]
\centering
\begin{tikzpicture}[x=1cm,y=1cm,font=\sffamily\scriptsize,
  wingtext/.style={align=left,inner sep=0pt},
  crate/.style={draw=black!75,fill=white,rounded corners=2pt,align=center,inner sep=3pt,line width=0.6pt},
  ext/.style={draw=black!60,fill=white,align=center,inner sep=3pt,line width=0.4pt},
  planned/.style={draw=black!50,dashed,fill=white,align=center,inner sep=3pt,line width=0.4pt},
  data/.style={draw=black!55,fill=yellow!14,align=center,inner sep=3pt,line width=0.4pt},
  zone/.style={draw=black!45,dashed,rounded corners=6pt,line width=0.5pt},
  zlabel/.style={font=\sffamily\scriptsize\itshape,text=black!70,inner sep=1pt},
  dep/.style={-{Stealth[length=4pt,width=3pt]},draw=black!55,line width=0.45pt},
  flow/.style={-{Stealth[length=5.5pt,width=4.5pt]},draw=blue!50!black,line width=1pt},
  num/.style={circle,fill=blue!50!black,text=white,inner sep=0.8pt,minimum size=9pt,font=\sffamily\bfseries\tiny}]

% ============================ (a) bow-tie ============================
\node[anchor=north west,font=\sffamily\footnotesize\bfseries,inner sep=0pt] at (0,0) {(a) One document, two readings};

\filldraw[fill=blue!7,draw=blue!40!black,line width=0.5pt] (0.2,-0.75) -- (7.0,-2.45) -- (7.0,-3.25) -- (0.2,-4.95) -- cycle;
\filldraw[fill=orange!12,draw=orange!55!black,line width=0.5pt] (17.4,-0.75) -- (10.6,-2.45) -- (10.6,-3.25) -- (17.4,-4.95) -- cycle;

\node[wingtext,anchor=west,text width=3.9cm] at (0.5,-2.85) {%
  {\footnotesize\bfseries Human reading}\\[1pt]
  the words the principal\\ reads and signs\\[3pt]
  {\ttfamily ALLOW git.pr.merge\\ ON repo "acme/*"\\ LIMIT 5 calls}};

\node[wingtext,align=right,anchor=east,text width=3.9cm] at (17.1,-2.85) {%
  {\footnotesize\bfseries Machine reading}\\[1pt]
  the typed tree the gate\\ evaluates\\[3pt]
  {\ttfamily allow: [git.pr.merge]\\ on: [(repo, "acme/*")]\\ limits: [Calls(5)]}};

\node[diamond,aspect=2.4,draw=black!80,fill=white,line width=0.7pt,align=center,inner sep=0pt,minimum width=3.6cm,minimum height=1.5cm] (knot) at (8.8,-2.85)
  {canonical text\\[-1pt]{\ttfamily id} = SHA-256};

\draw[flow] (6.0,-2.18) to[out=38,in=142] node[midway,above=1pt,text=black] {parse} (11.6,-2.18);
\draw[flow] (11.6,-3.52) to[out=218,in=322] node[midway,below=1pt,text=black] {render} (6.0,-3.52);

\node[data,text width=16.6cm,minimum height=0.85cm] (env) at (8.8,-5.75)
  {A signed envelope is the canonical text plus anchor identifier, assurance level, actor kind, signer key and validity window, under one signature.\\
   The identifier names the mandate in the {\ttfamily PARENT} line of every child, in quota counters and in receipts.};

% ============================ (b) crates and boundaries ============================
\node[anchor=north west,font=\sffamily\footnotesize\bfseries,inner sep=0pt] at (0,-6.75) {(b) Crates, data flow and trust boundaries};

% --- trusted inputs (top row)
\node[ext,text width=3.1cm,minimum height=0.95cm] (principal) at (1.9,-8.05) {{\bfseries principal}\\ signing device (trust anchor)};
\node[data,text width=7.2cm,minimum height=0.95cm] (config) at (9.0,-8.05) {{\bfseries launch configuration}\\ signed chain $\cdot$ actor key $\cdot$ tool policy map $\cdot$ clock};

% --- untrusted side
\node[ext,text width=3.1cm,minimum height=0.95cm] (model) at (1.9,-10.0) {{\bfseries language model}\\ prompts, injected text};
\node[ext,text width=3.1cm,minimum height=0.95cm] (harness) at (1.9,-11.75) {{\bfseries agent harness}\\ MCP client};
\node[planned,text width=3.1cm,minimum height=0.95cm] (hooks) at (1.9,-13.5) {harness hooks and plug-ins\\ (not in this release)};

% --- gate process
\node[crate,text width=3.3cm,minimum height=1.2cm] (mcp) at (7.0,-10.6) {{\ttfamily\bfseries handfest-mcp}\\ MCP stdio proxy};
\node[crate,text width=3.3cm,minimum height=1.2cm] (cli) at (11.0,-10.6) {{\ttfamily\bfseries handfest-cli}\\ fmt, keygen, sign,\\ verify, check, audit};
\node[crate,text width=3.3cm,minimum height=1.9cm] (anchors) at (7.0,-12.95) {{\ttfamily\bfseries handfest-anchors}\\ AnchorVerifier, AnchorRegistry\\ Ed25519: implemented\\ {\itshape passkey, wallet, QES: not written}};
\node[crate,text width=3.3cm,minimum height=1.9cm] (ledger) at (11.0,-12.95) {{\ttfamily\bfseries handfest-ledger}\\ ReceiptLedger\\ SqliteLedger: receipts\\ and quota counters};
\node[crate,fill=green!8,text width=6.6cm,minimum height=1.25cm] (core) at (8.6,-15.35) {{\ttfamily\bfseries handfest-core}\quad no\_std, no I/O, no clock\\ mandate type $\cdot$ parser and renderer $\cdot$ narrowing $\cdot$ evaluator};

% --- protected side
\node[ext,text width=3.0cm,minimum height=1.2cm] (upstream) at (15.75,-10.6) {{\bfseries upstream MCP server}\\ child process of the proxy};
\node[ext,text width=3.0cm,minimum height=0.95cm] (tools) at (15.75,-12.7) {{\bfseries tools and resources}\\ repositories, payments, \dots};

% --- audit side
\node[data,cylinder,shape border rotate=90,aspect=0.22,text width=2.5cm,minimum height=1.45cm] (file) at (12.55,-17.85) {{\bfseries ledger file}\\ SQLite,\\ rollback journal};
\node[ext,text width=3.0cm,minimum height=0.95cm] (auditor) at (15.75,-17.85) {{\bfseries offline auditor}\\ {\ttfamily handfest audit}};

% --- zones
\begin{scope}[on background layer]
  \node[zone,fill=red!4,fit={(model)(harness)(hooks)([yshift=-11pt]hooks.south)},inner sep=7pt] (zu) {};
  \node[zone,fill=black!3,fit={(config)(mcp)(cli)(anchors)(ledger)(core)([yshift=-11pt]core.south)},inner sep=8pt] (zt) {};
  \node[zone,fill=green!4,fit={(upstream)(tools)([yshift=-11pt]tools.south)},inner sep=7pt] (zp) {};
\end{scope}
\node[zlabel,anchor=south west] at ([xshift=5pt,yshift=3pt]zu.south west) {untrusted};
\node[zlabel,anchor=south west] at ([xshift=5pt,yshift=3pt]zt.south west) {gate process: trusted computing base};
\node[zlabel,anchor=south west] at ([xshift=5pt,yshift=3pt]zp.south west) {protected};

% --- dependencies (one way, downwards)
% both front ends depend on all three lower crates: a shared bus keeps the drawing readable
\coordinate (busL) at (mcp.south |- 0,-11.6);
\coordinate (busR) at (cli.south |- 0,-11.6);
\coordinate (busM) at (9.0,-11.6);
\draw[dep,-] (mcp.south) -- (busL) -- (busR) -- (cli.south);
\draw[dep] (busL) -- (anchors.north);
\draw[dep] (busR) -- (ledger.north);
\draw[dep] (busM) -- (busM |- core.north);
\draw[dep] (anchors.south) -- (anchors.south |- core.north);
\draw[dep] ([xshift=-14pt]ledger.south) -- ([xshift=-14pt]ledger.south |- core.north);

% --- data flow
\draw[flow] (principal.east) -- node[num,midway] {1} (config.west);
\draw[flow,black!55,line width=0.6pt] (model.south) -- (harness.north);
\draw[flow] ([yshift=4pt]harness.east) -- ++(1.15,0) |- node[num,pos=0.25] {2} ([yshift=5pt]mcp.west);
\draw[flow] ([yshift=-5pt]mcp.west) -- ++(-0.75,0) |- node[num,pos=0.25] {5} ([yshift=-6pt]harness.east);
\draw[flow] (config.south -| mcp.north) -- (mcp.north);
\draw[flow] ([xshift=22pt]mcp.north) -- ++(0,0.42) -| node[num,pos=0.25] {4} (upstream.north);
\draw[flow,black!55,line width=0.6pt] (upstream.south) -- (tools.north);
\draw[flow] ([xshift=44pt]ledger.south) -- node[num,pos=0.55] {3} ([xshift=44pt]ledger.south |- file.north);
\draw[flow] (file.east) -- node[num,midway] {6} (auditor.west);

% --- legend (drawn in the main picture; nested pictures would inherit node options)
\foreach \i/\y/\t in {
    1/-17.20/{the principal signs mandate text; the signed chain reaches the gate through its launcher},
    2/-17.62/{{\ttfamily tools/list} and {\ttfamily tools/call} as JSON-RPC over stdio},
    3/-18.04/{evaluate and record: receipt and counters in one transaction},
    4/-18.46/{the call is forwarded only after a permit has been committed},
    5/-18.88/{result, or error $-32001$ with reason and receipt hash},
    6/-19.30/{read-only audit: schema, hash chain, sequence mark, counters}}{
  \node[num] at (0.18,\y) {\i};
  \node[anchor=west,inner sep=0pt] at (0.48,\y) {\t};
}
\draw[dep] (0.02,-19.85) -- (0.6,-19.85);
\node[anchor=west,inner sep=0pt] at (0.72,-19.85) {crate dependency, one way};
\draw[flow] (4.6,-19.85) -- (5.18,-19.85);
\node[anchor=west,inner sep=0pt] at (5.3,-19.85) {data at run time};
\end{tikzpicture}
\caption{Handfest at a glance. (a)~The canonical text of a mandate is the single signed artifact; the tree the gate evaluates is obtained from it by a parser whose inverse is the renderer (Theorem~\ref{thm:bijection}), and its SHA-256 digest is the knot that ties both readings to every record. (b)~The five crates depend on each other in one direction only, toward the pure core. The gate runs outside the model: prompts, tool arguments and the bytes on the wire are untrusted, the launch configuration is trusted, and the upstream server is reached only after a permit has been committed to the ledger. Hooks and plug-ins for particular harnesses are not part of this release.}
\label{fig:arch}
\end{figure*}

% ---------------------------------------------------------------------------
% Figure: sequence diagram
% ---------------------------------------------------------------------------
\begin{figure*}[t]
\centering
\begin{tikzpicture}[x=1cm,y=1cm,font=\sffamily\scriptsize,
  actor/.style={draw=black!75,fill=black!6,rounded corners=2pt,align=center,inner sep=3pt,text width=2.3cm,minimum height=0.95cm,line width=0.6pt},
  life/.style={draw=black!40,dashed,line width=0.4pt},
  msg/.style={-{Stealth[length=5pt,width=4pt]},draw=black!85,line width=0.7pt},
  ret/.style={-{Stealth[length=5pt,width=4pt]},draw=black!85,dashed,line width=0.6pt},
  lab/.style={inner sep=1.5pt,fill=white,align=center,yshift=0.8pt},
  note/.style={draw=black!45,fill=yellow!14,align=left,inner sep=3pt,line width=0.4pt},
  act/.style={draw=black!70,fill=white,line width=0.5pt},
  tag/.style={font=\sffamily\scriptsize\bfseries,inner sep=2pt,fill=black!10,draw=black!55,line width=0.4pt}]

\foreach \x/\name/\title/\sub in {
    1.3/P/Human principal/reads and signs,
    4.3/A/Delegated agent/narrows and requests,
    7.3/G/Enforcement gate/MCP proxy or CLI,
    10.3/L/Receipt ledger/quotas and receipts,
    13.3/T/Protected tool/upstream server,
    16.3/U/Offline auditor/read-only}{
  \draw[life] (\x,-0.48) -- (\x,-14.3);
  \node[actor] (\name) at (\x,0) {{\bfseries \title}\\ \sub};
}

% activation bars
\filldraw[act] (7.19,-3.45) rectangle (7.41,-4.5);
\filldraw[act] (7.19,-5.4) rectangle (7.41,-12.45);
\filldraw[act] (10.19,-6.8) rectangle (10.41,-9.6);
\filldraw[act] (13.19,-10.5) rectangle (13.41,-11.0);
\filldraw[act] (10.19,-13.1) rectangle (10.41,-14.05);

% --- delegation
\node[note,anchor=west,text width=2.3cm] at (1.45,-1.05) {reads the canonical text and signs it};
\draw[msg] (1.3,-1.95) -- node[lab,above] {signed root mandate $M_0$} (4.3,-1.95);
\node[note,anchor=west,text width=2.3cm] at (4.45,-2.6) {narrows $M_0$ to a child $M_1$ and signs it};
\draw[msg] (4.3,-3.5) -- node[lab,above] {chain $M_0, M_1$} (7.19,-3.5);
\node[note,anchor=west,text width=2.35cm] at (7.55,-4.0) {verifies both signatures and $M_0 \sqsupseteq M_1$};

% --- per call
\draw[black!45,line width=0.4pt,dash pattern=on 1pt off 2pt] (0.05,-4.75) -- (17.55,-4.75);
\node[anchor=north west,font=\sffamily\scriptsize\itshape,text=black!75,inner sep=1.5pt,fill=white] at (0.05,-4.82) {for each tool call};

\draw[msg] (4.3,-5.45) -- node[lab,above] {tool call} (7.19,-5.45);
\node[note,anchor=west,text width=2.35cm] at (7.55,-6.0) {maps it to an action, a resource and facts};
\draw[msg] (7.41,-6.85) -- node[lab,above] {evaluate and record} (10.19,-6.85);
\node[note,anchor=north west,text width=6.35cm] at (10.6,-6.95) {%
  {\bfseries one transaction}\\
  read the counters and the last receipt hash\\
  evaluate (pure): proofs valid now, narrowing on every hop,\\
  \hspace*{0.6em}{\ttfamily WHERE} and {\ttfamily LIMIT} of every mandate in the chain\\
  append the receipt, permit or deny\\
  on a permit, charge every mandate in the chain\\
  commit};
\draw[ret] (10.19,-9.55) -- node[lab,above] {verdict, receipt hash} (7.41,-9.55);

% --- alternatives
\draw[black!55,line width=0.5pt] (2.95,-9.95) rectangle (14.65,-12.6);
\draw[black!55,line width=0.5pt,dashed] (2.95,-11.65) -- (14.65,-11.65);
\node[tag,anchor=north west] at (2.95,-9.95) {permit};
\node[tag,anchor=north west] at (2.95,-11.65) {deny};

\draw[msg] (7.41,-10.55) -- node[lab,above,pos=0.24] {forwards the call} (13.19,-10.55);
\draw[ret] (13.19,-11.0) -- node[lab,above,pos=0.76] {result} (7.41,-11.0);
\draw[ret] (7.19,-11.35) -- node[lab,above] {result} (4.3,-11.35);

\draw[ret] (7.19,-12.3) -- node[lab,above] {error $-32001$} (4.3,-12.3);
\node[anchor=west,align=left,inner sep=1.5pt,fill=white,font=\sffamily\scriptsize\itshape] at (7.55,-12.12) {with the reason and the receipt hash;\\ the tool is never contacted};

% --- audit
\draw[msg] (16.3,-13.15) -- node[lab,above,pos=0.3] {opens the ledger read-only} (10.41,-13.15);
\node[note,anchor=east,text width=5.3cm] at (16.15,-13.75) {replays the hash chain and checks the sequence mark and the counters; needs no write access};
\end{tikzpicture}
\caption{One delegation from signature to audit. The principal signs the text it has read; the delegate can only narrow it; the gate checks the chain once and then, for every call, has the ledger evaluate and record in a single transaction. The tool is contacted only after a permit has been committed. A denial is recorded too, and the tool never sees the call. The auditor needs the ledger file and nothing else.}
\label{fig:seq}
\end{figure*}

% ============================================================================
\section{Attacks and Hardening}
\label{sec:attacks}

The adversary we assume controls the model, and with it the arguments of every tool call and the bytes the harness sends to the gate. It may control a delegate that tries to widen what it was given, to pass as another kind of actor, or to split a budget. It may run a hostile or merely careless upstream server. Between runs it can write to the ledger file. We trust the principal's signing device, the gate binary, and whoever launches the gate, for the chain, the actor key, the policy map and the clock.

Table~\ref{tab:attacks} lists five families of attack. For each one the defense is in the code and a regression test exercises it.

\begin{table*}[t]
\centering
\caption{Attack families, the layer each targets, and where the regression tests live (paths relative to \texttt{crates/}).}
\label{tab:attacks}
\footnotesize
\setlength{\tabcolsep}{4pt}
\begin{tabular}{@{}L{2.6cm}L{1.7cm}L{4.5cm}L{4.6cm}L{3.2cm}@{}}
\toprule
Attack & Layer & What goes wrong & Defense & Tests \\
\midrule
Second preimage on the identifier & core &
A field containing a line feed renders to text that parses as different clauses. &
Private fields, validating constructors, no control characters in any string. &
\code{handfest-core/tests/core_tests.rs} \\
\addlinespace[2pt]
Negative usage; float rounding & core, both front ends &
A negative recorded total widens a budget; an amount just above a limit rounds down to it in binary floating point. &
Negative bounds, amounts and totals refused; number tokens kept as text and converted to exact decimals. &
\code{handfest-core/tests/gate_tests.rs}, \code{handfest-cli/tests/cli_tests.rs} \\
\addlinespace[2pt]
Unsigned envelope fields & anchors &
The validity window or the assurance level is changed without touching the signed text. &
One signature over a domain-separated, length-prefixed encoding of every field. &
\code{handfest-anchors/tests/anchors_tests.rs} \\
\addlinespace[2pt]
Framing and parser differentials & MCP proxy &
A bare carriage return, a duplicate or case-folded key, a response with both \texttt{result} and \texttt{error}, or a flood of requests makes proxy and server disagree. &
LF framing that rejects stray CR; duplicate, non-ASCII and case-colliding keys refused; \texttt{result} and \texttt{error} exclusive; bounded queues that fail closed. &
\code{handfest-mcp/tests/mcp_tests.rs}, \code{case_collision_tests.rs}, \code{session_transport_tests.rs} \\
\addlinespace[2pt]
Ledger file tampering & ledger &
The newest receipts are deleted, a trigger or a changed collation is planted, or an index is corrupted but not its table. &
Sequence high-water mark, exact comparison of the schema, full integrity check, reconciliation of counters. &
\code{handfest-ledger/tests/ledger_tests.rs}, \code{handfest-ledger/tests/hardening_and_denial_tests.rs} \\
\bottomrule
\end{tabular}
\end{table*}

\subsection{A Second Preimage Through the Renderer}
\label{sec:attack-preimage}

The identifier is a hash of rendered text, so two different trees that render to the same text would share an identifier and a signature. If trees could be built with arbitrary field values this would be easy. Give a resource selector the kind
\begin{lstlisting}
z_pad "x"\nWHERE resource.author ==
\end{lstlisting}
(with a real line feed for \texttt{\textbackslash n}) and the pattern \texttt{alice}. The renderer prints an \texttt{ON} line and then, on a line of its own, text that the parser reads as a \texttt{WHERE} clause. A tree with no conditions and a tree with one now have the same canonical text, and a signature on the restricted mandate also covers the unrestricted one.

This is why the type has private fields and only validating constructors, why kinds are identifiers and no string may contain a control character, and why Theorem~\ref{thm:bijection} is stated for the type and not for some larger set of trees.

\subsection{Numbers}
\label{sec:attack-numeric}

Two mistakes in handling numbers would each let a spender pass a limit. The first concerns state. If the recorded total of a sum quota is negative, because the ledger was corrupted or edited, then total plus amount stays under the bound for amounts far above it. The evaluator refuses a negative bound, a negative amount and a negative recorded total; the ledger refuses to load one; and the audit rejects it.

The second concerns input. Common JSON libraries read numbers into binary floating point, the precision that RFC~8259 tells interoperable software to expect~\cite{rfc8259}. Under \texttt{LIMIT 100.0 USD sum(resource.amount)}, the amount \texttt{100.000000000000001} is exactly $100$ as a 64-bit float and passes. Both front ends therefore read JSON with a parser in the core that keeps each number token as text, rejects duplicate keys at every depth, and converts tokens to exact decimals, with or without an exponent, refusing any that would need more than 28 fractional digits. The amount above is denied.

\subsection{Envelope Fields}
\label{sec:attack-envelope}

A signature over the mandate text alone leaves the rest of the envelope open. Anyone holding a signed mandate could extend its validity by years, or relabel a key file as a qualified signature. The signed payload is therefore a fixed domain prefix, then the anchor identifier, the assurance level, the actor kind and the signer key, each preceded by its length, then the two ends of the validity window, then the length-prefixed text. The verifier recomputes this payload from the envelope it was given, and decoding rejects unknown envelope fields, so nothing unsigned can ride along.

\subsection{The Wire}
\label{sec:attack-wire}

A proxy that parses messages differently from the server behind it can be made to approve one request and forward another. Four cases are handled explicitly.

\emph{Line framing.} MCP's stdio transport delimits messages by newlines~\cite{mcp2025spec}. If one side also takes a bare carriage return as the end of a line, a second message can hide inside the first. The proxy accepts LF and CRLF line ends, rejects any frame from either side that contains another carriage return, and refuses to emit one. Frames are limited to one mebibyte by default.

\emph{Key matching.} Go's \code{encoding/json} matches object keys to struct fields without regard to case, with the folding of \code{strings.EqualFold}~\cite{golang_json}. Under that folding the long~s (U+017F) equals \texttt{s} and the Kelvin sign (U+212A) equals \texttt{k}, so a key that an exact-match proxy does not recognize can still bind to a field in a Go server. Duplicate keys, whose treatment RFC~8259 calls unpredictable~\cite{rfc8259}, offer the same opening. The proxy's JSON reader rejects duplicates. Message envelopes and call parameters are closed sets of keys. In tool arguments every key must be ASCII, no two keys of an object may differ only in case, and no key may differ only in case from one the policy map names.

\emph{Result and error.} JSON-RPC requires a response to carry \texttt{result} or \texttt{error} and never both~\cite{jsonrpc2}. A proxy that looks only for the presence of \texttt{error} can be handed a tool list in \texttt{result} beside \texttt{"error": null}; it relays what it takes for an error, and the client reads the unfiltered list. The proxy rejects every response that has both members or neither, whatever their values.

\emph{Flooding.} The proxy keeps one request in flight, a queue of client frames waiting behind it, and a list of canceled requests whose responses are still owed. Each is bounded at 1{,}024 entries, as is the number of server messages that may arrive before a response. Evicting old entries to make room would let a client push out the record that marks a pending \texttt{tools/list} response for filtering. Exceeding any bound ends the session with an error instead.

\subsection{The Ledger File}
\label{sec:attack-ledger}

Someone who can write to the ledger file between runs can try four things that a plain replay of the hash chain misses.

\emph{Deleting the newest receipts.} The remaining prefix is a valid chain. The receipt table uses \texttt{AUTOINCREMENT}, for which SQLite records the largest sequence number ever issued in a table of its own~\cite{sqlite_autoinc}. The audit requires that number to equal the count of receipts, and every call counter to equal the number of permits that name its mandate.

\emph{Planting a trigger}, for instance one that zeroes the call counters whenever a receipt is inserted, so that the gate's own next write resets the quotas. When a ledger is opened for writing, and again on every audit, each row of SQLite's schema table is compared with the DDL the ledger expects, and anything else is rejected.

\emph{Changing a collation.} Rebuilding a counter table with a case-insensitive collation makes SQLite conflate units that the evaluator distinguishes. The same comparison catches it.

\emph{Corrupting an index but not its table.} A scan of the table still succeeds, and so does SQLite's quick check, which does not compare index content with table content~\cite{sqlite_pragma}. The ledger runs the full integrity check.

One hazard came from our own precautions. Before SQLite opens a file for an audit, the ledger reads the file header itself, to turn away files that are not rollback-journal databases without creating side files. On Unix, closing any descriptor of a file drops every POSIX lock the process holds on it~\cite{sqlite_corrupt}, including SQLite's. The ledger therefore keeps a per-process registry of open ledgers, keyed by device and inode, and skips its own inspection when the process already has the file open.

What a single file cannot do is prove that it is the right file. An attacker who rewrites receipts and recomputes every later hash, who restores an older copy of the database, or who substitutes a freshly built one leaves a file that is consistent with itself. Sum counters have a narrower gap of the same kind: receipts do not record the amounts charged, so replacing one well-formed total with another goes unnoticed. Detecting any of these needs a commitment kept somewhere else, such as the latest receipt hash held by a second party, or a transparency log~\cite{rfc6962}.

% ============================================================================
\section{Implementation and Status}
\label{sec:impl}

Table~\ref{tab:crates} gives the size of each crate at the revision this paper describes. The workspace is Rust~2021 with a minimum compiler version of~1.88. Unsafe code is forbidden by a workspace-wide lint. Clippy runs with warnings as errors and with limits on function length, nesting depth and argument count.

\begin{table}[t]
\centering
\caption{The workspace by crate: lines of Rust under \texttt{src/} and under \texttt{tests/}, and direct external dependencies. Test modules that live under \texttt{src/} are counted there.}
\label{tab:crates}
\footnotesize
\setlength{\tabcolsep}{3.5pt}
\begin{tabular}{@{}lrrL{3.25cm}@{}}
\toprule
Crate & \texttt{src/} & \texttt{tests/} & Depends on \\
\midrule
\code{handfest-core} & 5{,}693 & 3{,}133 & \code{sha2} \\
\code{handfest-anchors} & 686 & 998 & \code{ed25519-dalek}, \code{rand_core}, \code{hex}, \code{serde}, \code{serde_json} \\
\code{handfest-ledger} & 2{,}921 & 2{,}585 & \code{rusqlite} (bundled SQLite), \code{serde}, \code{serde_json} \\
\code{handfest-mcp} & 6{,}523 & 9{,}227 & \code{clap}, \code{serde}, \code{serde_json} \\
\code{handfest-cli} & 691 & 745 & \code{clap}, \code{serde}, \code{serde_json} \\
\midrule
Total & 16{,}514 & 16{,}688 & \\
\bottomrule
\end{tabular}
\end{table}

The tests are about as large as the code. They comprise unit tests, integration tests that drive both binaries as subprocesses, and adversarial tests for every attack in Section~\ref{sec:attacks}. Continuous integration runs the formatter, Clippy and the test suite, which are also the three commands needed to reproduce our checks:
\begin{lstlisting}
cargo fmt --all -- --check
cargo clippy --workspace --all-targets --locked \
  -- -D warnings
cargo test --workspace --locked
\end{lstlisting}
We do not measure coverage, we do not run mutation testing, and none of the formal statements in this paper has been mechanized: Theorems~\ref{thm:bijection} and~\ref{thm:sound} and Propositions~\ref{prop:decide} and~\ref{prop:budget} rest on the arguments given and on those tests.

We report no timings. What can be said without measuring is that narrowing has the bound of Proposition~\ref{prop:decide}, and that evaluating a request repeats the narrowing check on every hop and then evaluates every mandate. For a chain whose texts total $N$ bytes, a request that occupies $s \ge 1$ bytes with its facts, and a usage snapshot of $u$ entries, that takes time $O(N^2 + N s \log(N + s + u))$: each clause is matched against the request in time proportional to the product of their sizes, and the logarithm pays for the lookups of facts, set members and counters. A chain cannot be longer than the root's hop count plus one. Each decision is one SQLite transaction with full synchronous writes.

% ============================================================================
\section{Limitations}
\label{sec:limits}

\paragraph{One trust anchor.}
Only Ed25519 key files can sign. The assurance levels above a local key are defined and unreachable, and nothing in a mandate or in the gate can demand a minimum level. Passkeys, wallet signatures, selective-disclosure credentials and qualified signatures each need a verifier behind the existing trait.

\paragraph{No pinned roots and no revocation.}
A root mandate is self-asserted, and the party that launches a gate is trusted to supply the right chain. A broker that serves several principals needs a list of accepted roots. A mandate is valid until its envelope window or its own time conditions run out, and short windows are the only control today. Revocation could enter the pure evaluator as one more input, a signed status list for example, but it does not exist yet.

\paragraph{Facts need a source.}
In the proxy, fixed facts are the operator's word and lifted facts describe the call itself. A condition such as ``continuous integration is passing'' needs a source that attests to it at the time of the call. The evaluator accepts such facts; nothing yet produces them.

\paragraph{One ledger, one file.}
Proposition~\ref{prop:budget} holds for a single ledger. There is no distributed counter, and a ledger file cannot show that it has not been rolled back or replaced (Section~\ref{sec:attack-ledger}).

\paragraph{The screen and the law.}
The argument of Section~\ref{sec:wysiwys} assumes that the signing device displays the text it signs. We also make no claim that a mandate is a valid power of attorney in any jurisdiction; what it offers is evidence of what was signed, and its weight depends on the anchor and on the forum.

\paragraph{A small language and one transport.}
Mandates cannot relate two fields, do arithmetic, or negate a condition. The proxy speaks MCP over stdio only, serves one request at a time, and answers requests that a server addresses to the client with an error. There are no hooks for particular harnesses.

% ============================================================================
\section{Conclusion}
\label{sec:conclusion}

Authority handed to an agent should be something a person has read and signed and a machine can check. We have argued that the two requirements meet only if the signed bytes are the words themselves, and that this is what turns a signature into evidence of intent. A language small enough to guarantee it is also small enough for narrowing to be decided exactly and cheaply. Around that language, the rest of Handfest is conventional engineering done carefully: a pure evaluator, counters charged along the whole chain, a ledger that records refusals as well as permits, and a proxy that distrusts both ends of its pipe. The missing pieces, stronger anchors and revocation above all, fit behind interfaces that already exist.

\section*{AI Assistance Disclosure}
LLM-assisted tools were used throughout this project: in writing and reviewing the implementation, in generating tests, and in drafting, editing and typesetting this manuscript and its figures. No AI system is an author. The author, Valentin Yanakiev, takes full responsibility for every claim, proof, citation and line of code.

\section*{Artifact Availability}
Handfest is open source under the MIT or Apache-2.0 license, at the user's option, at \url{https://github.com/feedback-loop-ai/handfest}. This paper describes the code at commit \texttt{344df89}.

\bibliographystyle{plainnat}
\bibliography{references}

\end{document}
